\section{Definition of $m_2$}

\subsection{Concept}

$m_2$ denotes the domain of projection
in which relational states remain uniquely recoverable.

$m_2$ is not a space,
not a dimension,
and not a physical layer.

It is a validity domain
for physically interpretable statements.

\subsection{Formal definition}

For an explicit relational context $C$, let
\[
\Pi_C : \mathcal{U}_C \rightarrow \mathcal{M}_C
\]
be the corresponding projection. The context includes the applicable
TRD coupling, reference structure, and projection conditions.

Then
\[
m_2(C) := \{ u \in \mathcal{U}_C \mid \Pi_C \text{ is injective on } u \}
\]

Equivalently:
\[
u \in m_2(C)
\;\Longleftrightarrow\;
\exists!\,\Pi_C^{-1}(\Pi_C(u))
\]

Injectivity and uniqueness are always evaluated relative to this fixed context.
Different explicit contexts may yield different values without contradiction.
If the context is underdetermined, no unique physical back-reference exists
and the statement is outside $m_2(C)$.

\subsection{Properties}

Within $m_2(C)$:
\begin{itemize}
    \item physical quantities are well-defined,
    \item causal reconstruction is possible,
    \item comparisons between states are unambiguous.
\end{itemize}

Outside $m_2(C)$,
formal computation may continue,
but physical interpretation is no longer unique.

\subsection{Epistemic status}

$m_2$ is not fixed.

It expands or contracts
with available resolution,
measurement access,
and relational distinguishability.

$m_2$ represents
the current boundary
of reliable physical description.
